\chapter{Physical Commodity Markets}\label{ch:m3:physical-commodity-markets}

A tanker finishes loading 700\,000 barrels of North Sea crude on a Tuesday.
The buyer and the seller signed their contract weeks ago, and neither of them
yet knows the price. The contract says: the average of a published daily
assessment of North Sea crude over five business days, two before the date on
the loading document, that date, and two after, plus forty cents. Both parties
will learn the price the following Monday evening, when the last of the five
assessments is published. Until then the refinery that bought the oil holds a
cargo whose cost moves every day, and it has to decide whether, and how, to fix
it. This chapter describes the layer of the commodity world that such contracts
make up, the \omterm{def:m3:physical-commodity-markets:layers}{physical market}, and the tools by which traders turn its floating
prices into the fixed prices a business can plan with.

\section{Physical and paper}

Commodity markets have two layers. One moves goods: ships, pipelines, tanks,
silos and warehouses, with contracts that name a quality, a quantity, a place
and a date. The other moves only money: futures, swaps and options whose
settlement depends on a price. Books 1 and 2 of this series dealt almost
entirely with the second kind of instrument; in commodities the first kind comes
first, because it is where prices are made.

\begin{definition}[Physical market, paper market]\label{def:m3:physical-commodity-markets:layers}
The \emph{physical market}\index{physical market} of a commodity is the set of
contracts for the delivery of a specified quantity and quality of the commodity at
a specified place and time. Its \emph{paper market}\index{paper market} is the set
of contracts on the commodity's price that are normally settled in cash or closed
out before delivery: futures, swaps and options.
\end{definition}

The boundary is porous. A futures contract with physical delivery
(One Quant Book 1, chapter 21) is paper for almost all of its life and physical
at its end; an exchange for physical turns one into the other. The \omterm{def:m3:physical-commodity-markets:layers}{paper market}'s
own workhorse outside exchanges is the swap.

\begin{definition}[Commodity swap]\label{def:m3:physical-commodity-markets:swap}
A \emph{commodity swap}\index{commodity swap} exchanges, for each period of its
life, a fixed price for the average of a published floating price over the period,
both applied to an agreed volume, and is settled in cash for the difference.
\end{definition}

\begin{example}[A monthly crude swap]\label{ex:m3:physical-commodity-markets:swap}
An airline buys a swap on 100\,000 barrels a month for the next quarter at a fixed
\qty{82.00}{\$/\barrel} against the monthly average of a crude benchmark. If
the benchmark averages \qty{85.40}{\$/\barrel} in the first month, the airline
receives $3.40 \times 100\,000 = \$340\,000$; if it averages \qty{79.10}{\$/\barrel}
in the second, it pays $\$290\,000$. Whatever it pays for its fuel in the
\omterm{def:m3:physical-commodity-markets:layers}{physical market}, its net cost per barrel is fixed near \$82, up to the difference
between the fuel it burns and the benchmark it hedged, which is the subject of
the last section.
\end{example}

\section{Trading houses and the physical value chain}

Between the producer that lifts a commodity and the consumer that burns, smelts
or mills it stand firms that buy, store, blend, ship and resell.

\begin{definition}[Commodity trading house]\label{def:m3:physical-commodity-markets:house}
A \emph{commodity trading house}\index{commodity trading house} is a firm whose
business is to buy physical commodities and resell them, earning the difference
between their value in one place, time or form and their value in another, and
hedging the flat price with paper.
\end{definition}

A trading house transforms a commodity in three ways, each of which a price
difference rewards: in space (a cargo bought at a loading port and sold at a
discharge port), in time (oil stored when the forward price pays for the tank)
and in form (grades blended into a specification a refinery wants). Pirrong's
study of the industry describes trading firms in exactly these terms, as
performers of physical transformations, whose profits come from the spreads
between locations, dates and qualities rather than from the level of prices.
That is why a well-run trading house hedges almost all of its flat-price
exposure: the level of oil prices is a risk it carries, not the one it is paid
for.

\begin{dated}{2026-09}{How large the largest trader is}\label{dat:m3:physical-commodity-markets:scale}
Vitol, among the largest energy traders, reported deliveries of 8 million barrels a day
of crude oil and products in 2025 (7.2 million in 2024), on a turnover of USD~343
billion. World oil demand was about 104 million barrels a day; the International
Energy Agency forecast in May 2026 that it would contract in 2026, to 104 million
barrels a day, after shipments through the Strait of Hormuz were disrupted.
\end{dated}

\begin{remark}[Why the business is concentrated]\label{rem:m3:physical-commodity-markets:scale}
A physical trade ties up working capital for the whole voyage: a
700\,000-barrel cargo at \qty{80}{\$/\barrel} is USD~56 million, paid on
delivery and recovered weeks later. The firms that do this at scale are those
with the bank credit lines to finance it (\cref{ch:m3:getting-access-commodities-and-power})
and the flow to spread fixed costs over many cargoes.
\end{remark}

\section{Delivery terms and the physical contract}

A physical contract must say where the goods change hands, who pays to move
them, and who bears the risk of their loss on the way. International practice
uses standard terms for this.

\begin{definition}[Incoterms, free on board, cost insurance and freight]\label{def:m3:physical-commodity-markets:incoterms}
\emph{Incoterms}\index{Incoterms} are the standard delivery terms published by the
International Chamber of Commerce. Under \emph{free on board}\index{free on board}
(FOB) the seller delivers the goods on board the vessel nominated by the buyer at
the named port of shipment; risk passes there, and the buyer pays freight and
insurance. Under \emph{cost insurance and freight}\index{cost insurance and freight}
(CIF) risk also passes on board at the port of shipment, but the seller contracts
and pays for the carriage to the named destination port and for insurance against
the buyer's risk of loss.
\end{definition}

\begin{definition}[Bill of lading]\label{def:m3:physical-commodity-markets:bl}
A \emph{bill of lading}\index{bill of lading} is the document the carrier issues
on loading: a receipt for the goods, evidence of the contract of carriage, and a
document of title whose holder may claim the cargo at destination. Its date, the
B/L date, anchors most pricing formulas.
\end{definition}

\begin{omfigure}
\centering
\begin{tikzpicture}[font=\small,
  port/.style={draw,thick,rounded corners=2pt,align=center,minimum height=1.0cm,minimum width=2.6cm},
  arr/.style={-{Stealth[length=2.5mm]},thick}]
\node[port,draw=omDef,fill=omDef!8] (load) at (0,0) {loading port\\(seller's tank)};
\node[port,draw=omThm,fill=omThm!8] (disch) at (9.4,0) {discharge port\\(buyer's refinery)};
\draw[arr] (load.east) -- node[above=2pt] {voyage: freight, insurance, losses} (disch.west);
\draw[omProp,very thick] (2.05,-0.9) -- (2.05,0.9);
\node[omProp,font=\footnotesize\bfseries,align=center] at (2.05,1.25) {risk passes on board\\(FOB and CIF)};
\node[anchor=west,font=\footnotesize] at (-1.3,-1.35) {\textbf{FOB:} seller pays up to the ship's rail; buyer charters, insures, pays freight};
\node[anchor=west,font=\footnotesize] at (-1.3,-1.85) {\textbf{CIF:} seller also pays freight and insurance to the destination; risk still passes on board};
\node[anchor=west,font=\footnotesize] at (-1.3,-2.35) {\textbf{B/L date:} the day the bill of lading is issued; pricing periods are counted from it};
\end{tikzpicture}

\omcaption{The two delivery terms of seaborne commodity trade. Under both, the
buyer carries the risk of the voyage; under CIF the seller pays for it and prices
that cost into the cargo. Schematic.}\label{fig:m3:physical-commodity-markets:incoterms}
\end{omfigure}

A CIF price is therefore approximately the FOB price plus freight and insurance:
the two terms differ in who organises and pays for the voyage, not in who bears
its risk. Beyond the Incoterm, a physical contract names the grade and its
specification, the volume with an operational tolerance (a cargo is rarely
exactly the size agreed), the loading or delivery date range, the inspection
agent whose measurement of quantity and quality is binding, the payment terms
(often by letter of credit, thirty days after the B/L date; \cref{ch:m3:getting-access-commodities-and-power}) and the price
formula.

\section{Price-reporting agencies and assessment windows}

Most physical cargoes do not trade on an exchange, so there is no exchange price
to write into their formulas. The prices come from specialised publishers.

\begin{definition}[Price-reporting agency, assessment window]\label{def:m3:physical-commodity-markets:pra}
A \emph{price-reporting agency}\index{price-reporting agency} (PRA) is a publisher
that assesses and publishes, every business day, the prices of physical commodities
and related over-the-counter contracts from the bids, offers and trades it observes
and from reports by market participants. An \emph{assessment
window}\index{assessment window} is a fixed daily period during which the agency
collects bids, offers and trades that it then uses to set its assessment at the
window's closing time.
\end{definition}

The two largest agencies in oil are Platts, part of S\&P Global, and Argus. Their
assessments are referenced by physical contracts and by the paper contracts that
hedge them, and so work as benchmarks. In October 2012, at the request of the
G20, the International Organization of Securities Commissions published
principles for oil price-reporting agencies whose assessments are referenced by
derivatives contracts: governance, the quality and integrity of the methodology,
and the management of conflicts of interest.

\begin{dated}{2026-09}{The North Sea window}\label{dat:m3:physical-commodity-markets:moc}
Platts assesses Dated Brent from its market-on-close process: bids, offers and trades
reported during a thirty-minute window, 16:00 to 16:30 London time each business day,
with the assessment reflecting the value of the most competitive grade at 16:30.
Participants submit bids and offers for cargoes of specified size, loading dates and
grade; the agency publishes them and the trades. A parallel window runs in Singapore
for Middle-East crude.
\end{dated}

\section{Pricing formulas, netbacks and hedging a physical deal}

\begin{definition}[Pricing period, price differential]\label{def:m3:physical-commodity-markets:formula}
The \emph{pricing period}\index{pricing period} of a physical contract is the set of
days whose benchmark assessments are averaged to set its price: for example the five
business days around the B/L date (written 2-1-2), or all days of the delivery
month. The \emph{price differential}\index{price differential} is the fixed premium
or discount added to that average; it carries the value of the grade, location and
timing relative to the benchmark, and is what the parties actually negotiate.
\end{definition}

A trader who buys at one place and sells at another compares them through the
value at a single point.

\begin{definition}[Netback]\label{def:m3:physical-commodity-markets:netback}
The \emph{netback}\index{netback} of a sale is its delivered price less the costs of
getting the commodity there (freight, insurance, losses in transit, duties and
terminal fees), expressed per unit at the point of origin.
\end{definition}

\begin{example}[Choosing a destination by netback]\label{ex:m3:physical-commodity-markets:netback}
A cargo can be sold delivered in Rotterdam at \qty{81.30}{\$/\barrel} with
\qty{1.10}{\$/\barrel} of freight, or in Singapore at \qty{83.10}{\$/\barrel} with
\qty{3.05}{\$/\barrel}; insurance costs \qty{0.03}{\$/\barrel} either way. The
\omterm{def:m3:physical-commodity-markets:netback}{netbacks} are \qty{80.17}{\$/\barrel} and \qty{80.02}{\$/\barrel}: Rotterdam, although
Singapore pays more.
\end{example}

A physical formula leaves the parties exposed to the benchmark until the last
pricing day. They hedge that exposure with paper on the same benchmark, or on a
closely related one, and what the hedge cannot remove has a name.

\begin{definition}[Basis risk]\label{def:m3:physical-commodity-markets:basisrisk}
\emph{Basis risk}\index{basis risk} is the risk that the price of the position
being hedged and the price of the hedging instrument move by different amounts: the
risk that remains in a hedged position because the basis between them changes.
\end{definition}

\begin{method}[Hedging a cargo priced over a period]\label{met:m3:physical-commodity-markets:hedge}
When the price is fixed on $n$ pricing days, each day fixes $1/n$ of the volume.
(i) On the day the exposure is created (the contract is signed and its price will
float), take a futures position of the same size and opposite sign to the exposure
that pricing will create: a buyer of floating-price oil who wants a fixed cost buys
futures now. (ii) On each pricing day, close $1/n$ of the futures at that day's
settlement. (iii) The hedged price is then the futures price at inception plus the
average basis over the pricing days plus the differential.
\end{method}

\begin{proposition}[What a hedged cargo costs]\label{prop:m3:physical-commodity-markets:hedged}
Let $F_d$ be the futures price and $b_d$ the physical-minus-futures basis on day
$d$, and let the pricing days be $d_1,\dots,d_n$. A buyer hedged by
\cref{met:m3:physical-commodity-markets:hedge} from day $0$ pays per barrel
\[
\frac1n\sum_{i=1}^n (F_{d_i}+b_{d_i}) + D - \frac1n\sum_{i=1}^n (F_{d_i}-F_0)
\;=\; F_0 + \frac1n\sum_{i=1}^n b_{d_i} + D ,
\]
where $D$ is the differential. The flat price has gone; the average basis remains.
\end{proposition}

\begin{proof}
The formula price is the first average plus $D$; the futures bought at $F_0$ and
sold in equal parts at $F_{d_i}$ gain the second average. Subtract.
\end{proof}

\Cref{fig:m3:physical-commodity-markets:costcdf} simulates the refinery of the
hook: futures move by \qty{1.80}{\$/\barrel} a day and the basis by
\qty{0.12}{\$/\barrel} a day, both as random walks. Unhedged, the cost of the
cargo has a standard deviation of \qty{8.13}{\$/\barrel}, about USD~5.7 million on
the cargo; hedged, of \qty{0.54}{\$/\barrel}, which is exactly the standard deviation
of the average basis over the five pricing days.

\begin{omfigure}
\begin{tikzpicture}
\begin{axis}[width=11cm,height=5cm,
  xlabel={cost of the cargo, \$/bbl},ylabel={cumulative probability},
  tick label style={font=\small},label style={font=\small},
  xmin=58,xmax=102,ymin=0,ymax=1,grid=major,grid style={gray!25},
  legend columns=2,legend style={font=\footnotesize,at={(0.5,-0.3)},anchor=north,draw=none,fill=none,/tikz/every even column/.append style={column sep=8pt}},
  table/col sep=comma]
\addplot[omDef,very thick] table[x=unhedged,y=p]{figdata/markets-3/01-physical-commodity-markets/costcdf.csv};
\addplot[omThm,very thick] table[x=hedged,y=p]{figdata/markets-3/01-physical-commodity-markets/costcdf.csv};
\legend{unhedged,hedged with futures over the pricing period}
\end{axis}
\end{tikzpicture}

\omcaption{Distribution of the cost of a cargo priced 2-1-2 around a B/L date twenty
business days away, with futures at \qty{80}{\$/\barrel}, a basis of
$-\qty{0.30}{\$/\barrel}$ and a differential of $+\qty{0.40}{\$/\barrel}$: 20\,000
simulated paths. The hedge removes the flat price and leaves the basis.
Illustrative volatilities; data: the chapter's tutorial.}\label{fig:m3:physical-commodity-markets:costcdf}
\end{omfigure}

\omterm{def:m3:physical-commodity-markets:basisrisk}{Basis risk} is small when the hedge and the cargo reference the same benchmark and
large when they do not. The largest physical benchmarks, the seaborne North Sea
barrel and the inland US barrel (\cref{ch:m3:crude-oil}), move together most of
the time and apart when logistics break
(\cref{fig:m3:physical-commodity-markets:differential}). In 2011 inland US crude
production outgrew the pipelines out of the storage hub at Cushing, Oklahoma, and
the inland barrel fell to a monthly-average discount of about \$27 in September
2011; in 2026 the seaborne barrel rose above it by as much as \$17.66 in the
April monthly average while shipments through the Strait of Hormuz were disrupted.
A hedger who had sold one benchmark against a cargo priced on the other would have
lost or gained those amounts on top of the flat price.

\begin{omfigure}
\begin{tikzpicture}
\begin{axis}[width=11cm,height=4.8cm,
  ylabel={WTI minus Brent, \$/bbl},
  tick label style={font=\small},label style={font=\small},
  xmin=2000,xmax=2026.75,ymin=-32,ymax=10,ytick={-30,-20,-10,0,10},grid=major,grid style={gray!25},
  xtick={2000,2004,2008,2012,2016,2020,2024},xticklabel style={/pgf/number format/1000 sep={}},
  table/col sep=comma]
\addplot[omDef,thick] table[x=t,y=spread]{figdata/markets-3/01-physical-commodity-markets/differential.csv};
\draw[gray] (axis cs:2000,0) -- (axis cs:2026.75,0);
\end{axis}
\end{tikzpicture}

\omcaption{The inland US benchmark minus the seaborne North Sea benchmark, monthly
averages of daily spot prices, January 2000 to August 2026. The pipeline bottleneck of
2011 and the Gulf disruption of 2026 are the two deepest discounts. Data: FRED series
DCOILWTICO and DCOILBRENTEU (US Energy Information Administration).}\label{fig:m3:physical-commodity-markets:differential}
\end{omfigure}

\section{Tutorial: pricing and hedging a cargo}\label{tut:m3:physical-commodity-markets:hedge}

\begin{tutorial}
\textbf{Goal.} Price a cargo from its formula, hedge it by
\cref{met:m3:physical-commodity-markets:hedge}, and measure what the hedge leaves.
\textbf{End state:} \cref{fig:m3:physical-commodity-markets:costcdf} and the
standard deviations \qty{8.13}{} and \qty{0.54}{\$/\barrel}.
\begin{tutsteps}
\item \textbf{The legs and the hedge.} A leg is bought or sold, priced over a set of
days; each pricing day fixes an equal share and the hedge trades that share.
\omcode{code/firm/physdeal/firm_physdeal.py}{28}{58}{Legs, exposure and the futures trades that offset it.}\label{lst:m3:physical-commodity-markets:legs}
\item \textbf{Simulate.} Futures and basis as random walks; cost unhedged and hedged,
and the theoretical hedged standard deviation from the covariance of a random walk.
\omcode{code/markets-3/01-physical-commodity-markets/python/m3_physical.py}{22}{39}{The refinery's cargo on 20\,000 paths.}\label{lst:m3:physical-commodity-markets:sim}
\item \textbf{Run} \texttt{m3\_physical.simulate\_cargo()},
\texttt{hedged\_sd\_theory()} and \texttt{fig\_physical.py}.
\end{tutsteps}
\textbf{What to change next.} Hedge on the bill-of-lading date only, instead of
over the \omterm{def:m3:physical-commodity-markets:formula}{pricing period}, and measure the residual; replace the basis random walk by a
mean-reverting one and see how much of the hedged risk remains.
\end{tutorial}

\section{Build: the physical deal model}\label{bld:m3:physical-commodity-markets:physdeal}

\begin{build}
\bfield{Purpose} Every physical trade of the miniature firm enters its books through
this model: what the price will be, which days still float, what freight and
insurance the firm pays, and which futures to trade each day to keep the book flat.
\bfield{Interface} \texttt{pricing\_days(event\_day, before, after)};
\texttt{formula\_price(assessments, days, differential)}; \texttt{Leg(side, volume,
days, differential)}; \texttt{exposure(legs, day)}; \texttt{hedge\_trades(legs)};
\texttt{cost\_split(incoterm, freight, insurance)}; \texttt{netback};
\texttt{locked\_margin}.
\bfield{Rules} Days are business-day indices; a missing assessment is an error, never
an interpolation; each pricing day fixes an equal share; futures lots are whole
(1\,000 barrels) or the schedule is rejected; FOB and CIF only.
\bfield{Acceptance tests} \texttt{code/firm/physdeal/tests/}: a 2-1-2 period; the
exposure of a back-to-back cargo is zero before and after both \omterm{def:m3:physical-commodity-markets:formula}{pricing periods} and
the full volume between them; hedge plus exposure is zero on every day; \omterm{def:m3:physical-commodity-markets:netback}{netback} and
the Incoterm split.
\bfield{Stretch} Post the futures trades to the position keeper of One Quant Book 1,
chapter 7; price-date calendars with holidays; volume tolerance at the buyer's
option.
\end{build}

\begin{omsources}
\item International Chamber of Commerce, \emph{Incoterms 2020}; ICC Academy, ``Place
of delivery and risk transfer in international trade contracts''.
\item IOSCO, \emph{Principles for Oil Price Reporting Agencies}, FR06/12, 5 October 2012
(Internet Archive copy), and the implementation report of 2013.
\item S\&P Global Platts, Dated Brent methodology and the market-on-close process.
\item C.~Pirrong, \emph{The Economics of Commodity Trading Firms}, 2014.
\item Vitol, \emph{2025 volumes and review}; IEA, \emph{Oil Market Report}, May and
August 2026.
\item US Energy Information Administration, \emph{Today in Energy}, ``Spread narrows
between Brent and WTI crude oil benchmark prices'', 2013; FRED series DCOILWTICO and
DCOILBRENTEU.
\end{omsources}

\section{Exercises}

\begin{exercise}[$\star$]\label{exo:m3:physical-commodity-markets:1}
A cargo is priced at the average of five assessments plus \qty{0.35}{\$/\barrel}.
The assessments are 81.20, 80.70, 79.90, 80.40 and 81.05. What is the price?
\end{exercise}

\begin{exercise}[$\star$]\label{exo:m3:physical-commodity-markets:2}
A seller offers a cargo at \qty{78.40}{\$/\barrel} FOB. Freight to the buyer's port
is \qty{1.35}{\$/\barrel} and insurance \qty{0.04}{\$/\barrel}. What CIF price is
equivalent, and who bears the risk of the voyage under each term?
\end{exercise}

\begin{exercise}[$\star$]\label{exo:m3:physical-commodity-markets:3}
A cargo sells delivered at \qty{84.10}{\$/\barrel}; freight is
\qty{2.20}{\$/\barrel}, insurance \qty{0.05}{\$/\barrel}, and 0.2\% of the volume
is lost in transit. Give the \omterm{def:m3:physical-commodity-markets:netback}{netback} per barrel loaded.
\end{exercise}

\begin{exercise}[$\star\star$]\label{exo:m3:physical-commodity-markets:4}
A refinery buys 700\,000 barrels priced 2-1-2 around the B/L date. How many
1\,000-barrel futures does it buy today, and how many does it sell on each pricing
day? What changes if one of the five days is a holiday on which no assessment is
published and the contract drops it?
\end{exercise}

\begin{exercise}[$\star\star$]\label{exo:m3:physical-commodity-markets:5}
In the simulation of \cref{fig:m3:physical-commodity-markets:costcdf} the hedged
standard deviation is \qty{0.54}{\$/\barrel}. What would it be if the daily
volatility of the basis doubled, and why does the futures volatility not enter?
\end{exercise}

\begin{exercise}[$\star\star$]\label{exo:m3:physical-commodity-markets:6}
Using \cref{ex:m3:physical-commodity-markets:netback}, by how much would Singapore's
delivered price have to rise for the trader to switch destinations?
\end{exercise}

\begin{exercise}[$\star\star\star$]\label{exo:m3:physical-commodity-markets:7}
\emph{Coding.} With \texttt{load\_differential}, find the month of the deepest WTI
discount to Brent since 2000 and its size, and count the months since January 2011
in which WTI averaged above Brent.
\end{exercise}

\begin{exercise}[$\star\star\star$]\label{exo:m3:physical-commodity-markets:8}
\emph{Find the flaw.} ``Our cargo prices around the B/L date, so we hedge it by
buying 700 futures on the B/L date and selling them once the price is known.''
Correct it.
\end{exercise}

\section{Problem: The Cargo That Priced Late}

\begin{problem}[{Weekend problem --- a back-to-back cargo with two pricing periods}]\label{pb:m3:physical-commodity-markets:1}
A trading house buys 700\,000 barrels FOB, priced 2-1-2 around the B/L date (day
10) at Dated plus \qty{0.20}{\$/\barrel}, and sells the same cargo CIF Rotterdam,
priced 2-1-2 around the discharge date (day 30) at Dated plus
\qty{1.10}{\$/\barrel}. It pays freight of \qty{0.65}{\$/\barrel} and insurance of
\qty{0.03}{\$/\barrel}. Dated moves by about \qty{1.80}{\$/\barrel} a business day.

\medskip\noindent\textbf{Part I --- The exposure.}
\begin{enumerate}
\item What is the trader's flat-price exposure before day 8?
\item How does it change on days 8 to 12?
\item What is it between days 13 and 27, and for how many days is it at its maximum?
\item How does it change on days 28 to 32?
\item Sketch the exposure against time.
\end{enumerate}

\medskip\noindent\textbf{Part II --- The hedge.}
\begin{enumerate}[resume]
\item What futures does the trader trade on each of days 8 to 12?
\item And on days 28 to 32?
\item What is the net futures position on day 20?
\item Why is no futures position needed before day 8?
\item Which instrument would you use if the futures did not settle on Dated?
\end{enumerate}

\medskip\noindent\textbf{Part III --- The money.}
\begin{enumerate}[resume]
\item Give the margin per barrel locked in by the hedge.
\item Give it for the cargo, in dollars.
\item Estimate the standard deviation of the unhedged P\&L.
\item What risks remain after the hedge?
\item How does a delay of the discharge by five days change the exposure and the hedge?
\end{enumerate}

\medskip\noindent\textbf{Part IV --- Judgement.}
\begin{enumerate}[resume]
\item Why might the trader choose a CIF sale rather than FOB?
\item What does the trader's credit line have to finance?
\item Why do trading houses call this business ``spread'' trading?
\item State the \emph{named result}: the locked margin of the cargo against its unhedged risk.
\item In one sentence: what does a physical trader earn?
\end{enumerate}
\end{problem}

\medskip
\begin{omfigure}
\begin{tikzpicture}
\begin{axis}[width=11cm,height=4.6cm,
  xlabel={business day},ylabel={thousand barrels / lots},
  tick label style={font=\small},label style={font=\small},
  xmin=0,xmax=35,ymin=-800,ymax=800,grid=major,grid style={gray!25},
  legend columns=2,legend style={font=\footnotesize,at={(0.5,-0.33)},anchor=north,draw=none,fill=none,/tikz/every even column/.append style={column sep=8pt}},
  table/col sep=comma]
\addplot[omDef,very thick,const plot] table[x=day,y=exposure]{figdata/markets-3/01-physical-commodity-markets/exposure.csv};
\addplot[omThm,very thick,dashed,const plot] table[x=day,y=hedge]{figdata/markets-3/01-physical-commodity-markets/exposure.csv};
\legend{physical exposure (thousand bbl),futures held (lots of 1\,000 bbl)}
\end{axis}
\end{tikzpicture}

\omcaption{The weekend problem's back-to-back cargo: the flat-price exposure builds
as the purchase prices in (days 8--12) and unwinds as the sale prices in (days
28--32); the futures position mirrors it. Data: the chapter's
code.}\label{fig:m3:physical-commodity-markets:exposure}
\end{omfigure}

\section{Interview questions}

\begin{interviewq}[$\star$ \iqroles{trader}]\label{iq:m3:physical-commodity-markets:1}
What is the difference between FOB and CIF, and which carries the voyage risk?
\end{interviewq}

\begin{interviewq}[$\star$ \iqroles{trader, researcher}]\label{iq:m3:physical-commodity-markets:2}
Why are most physical cargoes priced as an average over several days rather than at
a single price?
\end{interviewq}

\begin{interviewq}[$\star\star$ \iqroles{trader, risk}]\label{iq:m3:physical-commodity-markets:3}
You bought a cargo priced on the average of next month. When do you hedge, and how
much each day?
\end{interviewq}

\begin{interviewq}[$\star\star$ \iqroles{researcher}]\label{iq:m3:physical-commodity-markets:4}
A \omterm{def:m3:physical-commodity-markets:pra}{price-reporting agency}'s assessment and an exchange's settlement price are both
``the price''. How do they differ as data?
\end{interviewq}

\begin{interviewq}[$\star\star$ \iqroles{risk}]\label{iq:m3:physical-commodity-markets:5}
A trading house says it hedges all its flat-price risk. Where can it still lose
money?
\end{interviewq}

\begin{interviewq}[$\star\star\star$ \iqroles{developer, trader}]\label{iq:m3:physical-commodity-markets:6}
Design the service that tells a physical desk, every evening, its exposure by
benchmark and pricing day and the futures it should hold.
\end{interviewq}
